\documentclass{article}
\usepackage{amsmath}

\begin{document}

\section*{Two-asset minimum-variance portfolio}

Let two risky assets have variance parameters \(\sigma_1^2\) and \(\sigma_2^2\), and covariance
parameter \(\sigma_{12}\). Consider a portfolio that invests weight \(w \in \mathbb{R}\) in
asset~1 and weight \(1-w\) in asset~2. Its variance is
\[
V(w)
=
w^2 \sigma_1^2
+
(1-w)^2 \sigma_2^2
+
2 w (1-w) \sigma_{12}.
\]

Assume the quadratic coefficient is positive:
\[
\sigma_1^2 + \sigma_2^2 - 2\sigma_{12} > 0.
\]

Define
\[
w^\ast
=
\frac{\sigma_2^2 - \sigma_{12}}
{\sigma_1^2 + \sigma_2^2 - 2\sigma_{12}}.
\]

Then \(w^\ast\) minimizes \(V\) over all real weights \(w \in \mathbb{R}\) (this is a global
minimum over the whole real line, not a long-only \(0 \le w \le 1\) constraint). Equivalently,
for every \(w \in \mathbb{R}\),
\[
V(w^\ast) \leq V(w).
\]

No separate positivity of \(\sigma_1^2\) or \(\sigma_2^2\) is required; the only hypothesis is
that the quadratic coefficient \(\sigma_1^2 + \sigma_2^2 - 2\sigma_{12}\) is positive. The claim
is non-strict global minimality and asserts no uniqueness.

\end{document}
